\section{Physical solutions realizing local transport frames}
\label{sec:physical-amplitudes}

We now construct exact elasticity solutions with controlled leading
behavior for both the displacement and its actual divergence. We begin
by identifying the equations obeyed by these leading terms.
Fix the null vector $\theta$ from \eqref{eq:phase-notation}, and
write $D=D_\theta$ when convenient. Since
$\theta\cdot\theta=0$, conjugating the operators in
\eqref{eq:augmented-operators} gives
\[
 E_\tau^{-1}R(E_\tau a,E_\tau b)=R(a,b)+\tau R^\theta(a,b),
 \qquad
 E_\tau^{-1}S(E_\tau a,E_\tau b)=S(a,b)+\tau S^\theta(a,b),
\]
where
\begin{equation}
\label{eq:leading-operators}
\begin{aligned}
 R^\theta(a,b)
   &=2mDa+(Dm)a+\theta(kb+\nabla m\cdot a),\\
 S^\theta(a,b)
   &=2\ell Db+2(Dk)b+4\nabla m\cdot Da
                       +2(D\nabla m)\cdot a.
\end{aligned}
\end{equation}
For smooth vector amplitudes $a_0,a_1$, the first two terms of a trial
displacement give
\[
 E_\tau^{-1}\diverg\bigl(E_\tau(a_0+\tau^{-1}a_1)\bigr)
 =\tau\theta\cdot a_0+
   \bigl(\diverg a_0+\theta\cdot a_1\bigr)
   +\tau^{-1}\diverg a_1.
\]
An $O(1)$ leading divergence therefore requires $\theta\cdot a_0=0$,
while its scalar term also contains $\theta\cdot a_1$ and need not equal
$\diverg a_0$. We retain that scalar as the component $b$ of the leading
pair.
Define the fixed complex vector spaces
\begin{equation}
\label{eq:constrained-fibers}
 H_\theta=\{a\in\C^3:\theta\cdot a=0\},
 \qquad \Ecal_\theta=H_\theta\oplus\C.
\end{equation}
The transport equation on $\Ecal_\theta$ is
\begin{equation}
\label{eq:constrained-transport}
 R^\theta(a,b)=0,\qquad S^\theta(a,b)=0,
 \qquad (a,b)\in\Ecal_\theta.
\end{equation}
Indeed, putting $P^\theta=2mD+(Dm)$, we have
\begin{equation}
\label{eq:normal-transport-component}
 \theta\cdot R^\theta(a,b)=P^\theta(\theta\cdot a).
\end{equation}
Thus $R^\theta$ takes values in $H_\theta$ when $a$ does. On
$\Ecal_\theta$, the coefficient multiplying $D(a,b)$ in the two
transport equations is
\[
 \begin{pmatrix}
  2m\Id_{H_\theta}&0\\
  4\nabla m\cdot&2\ell
 \end{pmatrix}.
\]
It is invertible, so \eqref{eq:constrained-transport} is a square
rank-three system with principal part $D$.

For clarity, write $\theta=\alpha+i\beta$ with
$\alpha,\beta\in\R^3$. The null condition says
$\abs\alpha=\abs\beta=\rho>0$ and $\alpha\cdot\beta=0$.
Let $(y_1,y_2,t)$ be positively oriented orthonormal coordinates in
the directions $\alpha/\rho$, $\beta/\rho$, and their cross product.
Then, for $z=y_1+iy_2$,
\[
 D=\rho(\partial_{y_1}+i\partial_{y_2})
    =2\rho\partial_{\overline z}.
\]
We denote the planar disk $\{y\in\R^2:\abs y<R\}$ by $D_R^2$.

For each chosen transverse coordinate $t_0$, we will select a smooth
frame of \eqref{eq:constrained-transport} for $t$ near $t_0$ and then
realize its three columns by physical solutions.
Earlier systems complex geometric optics constructions also use
invertible planar transport frames \cite[Section~1]{EskinRalston2004}.
The next lemma records the local parameter dependence and inhomogeneous
solvability with transverse support needed here.

\begin{lemma}[Planar frames and inhomogeneous transports]
\label{lem:planar-frames}
Let $I\subset\R$ be an open interval and let
$A\in C^\infty(D_R^2\times I;\C^{n\times n})$.
For each fixed $t\in I$, the equation
\begin{equation}
\label{eq:planar-frame-equation}
 \partial_{\overline z}V=A(z,t)V
\end{equation}
has a smooth invertible matrix solution on $D_R^2$.
For every $R'<R$ and $t_0\in I$, there are an interval
$I_0\Subset I$ containing $t_0$ and a smooth invertible solution
$V$ on $D_{R'}^2\times I_0$.

Given this $V$, a radius $R''<R'$, and
$F\in C^\infty(D_{R'}^2\times I_0;\C^n)$, the equation
\begin{equation}
\label{eq:planar-inhomogeneous}
 \partial_{\overline z}v=Av+F
\end{equation}
has a smooth solution on $D_{R''}^2\times I_0$. If
$F(\,\cdot\,,t)=0$ for $t\notin K$, where $K\Subset I_0$ is
compact, the solution can be chosen with the same transverse
support property. These assertions also hold with $D$ in place of
$\partial_{\overline z}$, and for systems obtained by multiplying
$D$ by an invertible smooth matrix.
\end{lemma}

\begin{proof}
We first construct frames locally at a fixed parameter. For a
bounded function $f$ supported in a disk of radius $\delta$, the
Cauchy transform
\[
 (\mathcal Tf)(z)=\frac1\pi
       \int_{\C}\frac{f(w)}{z-w}\,\dd y_1(w)\,\dd y_2(w)
\]
satisfies $\partial_{\overline z}\mathcal Tf=f$ in distributions
and $\norm{\mathcal Tf}_{L^\infty}\le C\delta\norm f_{L^\infty}$.
The latter estimate follows by integrating $\abs{z-w}^{-1}$ over
the support disk. Choose a smooth cutoff $\chi$ supported in a
sufficiently small disk and equal to one on a smaller disk. On
bounded matrix functions, the equation
\[
 V=\Id_n+\mathcal T(\chi A V)
\]
is then solved by a convergent Neumann series. Shrinking the disk
until the operator norm is less than $1/3$ gives
$\norm{V-\Id_n}_{L^\infty}<1/2$, hence invertibility. The
distributional equation and interior elliptic regularity for
$\partial_{\overline z}$ show that $V$ is smooth on the smaller
disk. This proves local smooth invertible solvability.

If $V_i,V_j$ are two such local frames, then
$\partial_{\overline z}(V_i^{-1}V_j)=0$ on their overlap.
Their transition functions therefore define a holomorphic vector
bundle on $D_R^2$, whose underlying smooth bundle is trivial.
Its holomorphic frame bundle has fiber $\mathrm{GL}_n(\C)$.
The disk is Stein and contractible, so a continuous section of
this frame bundle exists. The Oka principle for sections of holomorphic
fiber bundles with complex homogeneous fibers, in the form of
\cite[Theorem~1.1, p.~1018]{Forstneric2009}, supplies a holomorphic
section; this homogeneous-fiber case goes back to Grauert
\cite{Grauert1958}. In the original smooth trivialization it is
a smooth invertible matrix solving
\eqref{eq:planar-frame-equation} throughout $D_R^2$.

To obtain local parameter dependence, choose such a frame $V_0(z)$
at $t_0$. After writing $V=V_0W$, the equation becomes
\[
 \partial_{\overline z}W=B(z,t)W,
 \qquad
 B(z,t)=V_0(z)^{-1}\bigl(A(z,t)-A(z,t_0)\bigr)V_0(z).
\]
Choose a cutoff $\chi\in C_c^\infty(D_R^2)$ equal to one near
$\overline D_{R'}^2$. On its compact support, $B(\,\cdot\,,t)$
tends uniformly to zero as $t\to t_0$. On a sufficiently small
interval $I_0\Subset I$, solve
\[
 W=\Id_n+\mathcal T(\chi B(\,\cdot\,,t)W)
\]
by a Neumann series with norm less than $1/3$. This gives an
invertible solution on $D_{R'}^2\times I_0$. The integral operator
depends smoothly on $t$ in operator norm on $L^\infty$; its
inverse does also. Hence all parameter derivatives of $W$ exist
in $L^\infty$. Differentiating its equation and applying interior
elliptic regularity successively gives smoothness jointly in
$(z,t)$, with bounds on compact subsets. Thus $V=V_0W$ has the
claimed parameter regularity.

Finally, choose $\chi\in C_c^\infty(D_{R'}^2)$ equal to one near
$\overline D_{R''}^2$. The formula
\[
 v=V\mathcal T(\chi V^{-1}F)
\]
solves \eqref{eq:planar-inhomogeneous} on the smaller disk.
The transform acts only in the planar variable, so it preserves
the specified support in $t$. Dividing a system with invertible
derivative coefficient by that coefficient, and using
$D=2\rho\partial_{\overline z}$, proves the final assertion.
\end{proof}

To carry a leading frame through physical transfer, we need control of
both the displacement--divergence pair and its derivatives.
Corollary~\ref{cor:physical-solvability} gives physical solvability with
an $L^2$ estimate independent of $\tau$. We will make the
correction small enough after taking a divergence by first constructing
a compatible expansion of truncation order $N\ge4$. The proof below
obtains a physical residual of size $O(\tau^{1-N})$ and a conjugated
divergence correction of size $O(\tau^{2-N})$. The resulting strong
convergence identifies the first medium's leading columns and, by exterior
agreement, the exterior values of the transferred columns. The weighted
$O(\tau)$ derivative bound controls the first-order forcing obtained by
applying the difference of the two normalized augmented operators to
the first physical pair.

\begin{proposition}[Physical realization of local transport frames]
\label{prop:physical-amplitudes}
For every nonzero null vector $\theta$ and every
$t_0\in(-2r,2r)$, there exist an open interval
$J\subset(-2r,2r)$ containing $t_0$ and three smooth solutions
$(a_0^{\nu},b_0^{\nu})$, $1\le\nu\le3$, of
\eqref{eq:constrained-transport} on a neighborhood of
$\overline B_{5r}$, with the following properties.

The three columns form a frame of $\Ecal_\theta$ at every point
of $D_{2r}^2\times J$. For each $\nu$ and all sufficiently large
$\tau$, there is $u_\tau^{\nu}\in C^\infty(B_{5r};\C^3)$ with
$Lu_\tau^{\nu}=0$ such that, on putting
$U_\tau^{\nu}=(u_\tau^{\nu},\diverg u_\tau^{\nu})$,
\begin{equation}
\label{eq:physical-leading-limit}
\begin{aligned}
 E_\tau^{-1}U_\tau^{\nu}
   &\longrightarrow (a_0^{\nu},b_0^{\nu})
                  &&\text{in }L^2(B_{4r}),\\
 \norm{E_\tau^{-1}\nabla U_\tau^{\nu}}_{L^2(B_{4r})}
   &=O(\tau).
\end{aligned}
\end{equation}
\end{proposition}

\begin{proof}
We construct compatible formal amplitudes and then correct the
truncated displacement using the physical operator.
Fix an integer $N\ge4$ and radii
$R_0>R_1>\cdots>R_N>5r$. The coefficients are smooth on all of
$\R^3$. Apply Lemma~\ref{lem:planar-frames} to the rank-three
system \eqref{eq:constrained-transport}, starting on a disk
strictly larger than $D_{R_0}^2$. It supplies a smooth
invertible frame
\[
 \mathcal F(y,t):\C^3\longrightarrow\Ecal_\theta
 \quad\text{on }D_{R_0}^2\times I_0,
\]
where $I_0\Subset(-2r,2r)$ contains $t_0$. Choose an interval
$J$ containing $t_0$ with $\overline J\subset I_0$, and choose
$\chi\in C_c^\infty(I_0)$ equal to one on $J$. Multiply each
column of $\mathcal F$ by $\chi(t)$ and extend it by zero in
$t$ to obtain $(a_0^{\nu},b_0^{\nu})$ on
$D_{R_0}^2\times\R$. Since $D\chi=0$, these columns still
solve \eqref{eq:constrained-transport}. They are a frame on
$D_{2r}^2\times J$, and all have transverse support in the
fixed compact set $K=\supp\chi\Subset I_0$.

We construct a physical solution for any one of these columns
and suppress $\nu$. Higher-order transport expansions for systems appear
in Eskin and Ralston
\cite[Section~2, equations~(2.2)--(2.4)]{EskinRalston2004}.
The recursion here must cancel the vector residual
while making the conjugated actual divergence have leading term $b_0$.
Its scalar equation makes the next normal constraint compatible with
the vector equation;
the final correction will be made only in the physical displacement
equation. For $1\le j\le N$, we seek smooth coefficients $(a_j,b_j)$
on successively smaller disks such that
\begin{equation}
\label{eq:physical-recursion}
\begin{aligned}
 R^\theta(a_j,b_j)&=-R(a_{j-1},b_{j-1}),\\
 S^\theta(a_j,b_j)&=-S(a_{j-1},b_{j-1}),\\
 \theta\cdot a_j&=-(\diverg a_{j-1}-b_{j-1}).
\end{aligned}
\end{equation}
Throughout this recursion, write $q_j=\diverg a_j-b_j$ and set
$a_{-1}=b_{-1}=q_{-1}=0$.
The last equation in \eqref{eq:physical-recursion} is the
displacement--divergence constraint at order $j$.

We first verify the compatibility of these three equations.
Conjugating $P$ gives $E_\tau^{-1}PE_\tau=P+\tau P^\theta$.
Comparing coefficients of $\tau$ in the conjugated identity
\eqref{eq:augmented-identity} yields, for arbitrary $(a,b)$,
\begin{equation}
\label{eq:transport-compatibility}
 \diverg R^\theta(a,b)+\theta\cdot R(a,b)-S^\theta(a,b)
 =P(\theta\cdot a)+P^\theta(\diverg a-b).
\end{equation}
Suppose the preceding recursive equations have been solved.
Apply \eqref{eq:transport-compatibility} to $(a_{j-1},b_{j-1})$.
For $j\ge2$, the preceding vector and scalar equations, followed by
\eqref{eq:augmented-identity}, give
\[
 \begin{aligned}
 \diverg R^\theta(a_{j-1},b_{j-1})
       -S^\theta(a_{j-1},b_{j-1})
   &=-\diverg R(a_{j-2},b_{j-2})+S(a_{j-2},b_{j-2})\\
   &=-Pq_{j-2}.
 \end{aligned}
\]
The preceding normal constraint gives the same term on the other side:
\[
 P(\theta\cdot a_{j-1})=-Pq_{j-2}.
\]
For $j=1$, both expressions are zero by the homogeneous leading
equations and $q_{-1}=0$. In every case these two terms cancel in
\eqref{eq:transport-compatibility}, leaving
\begin{equation}
\label{eq:recursion-normal-compatibility}
 \theta\cdot R(a_{j-1},b_{j-1})=P^\theta q_{j-1}.
\end{equation}

Choose a constant vector $n_\theta\in\C^3$ with
$\theta\cdot n_\theta=1$, and set
$a_j^{\mathrm p}=-q_{j-1}n_\theta$, $b_j^{\mathrm p}=0$.
After subtracting this particular pair from the desired
$(a_j,b_j)$, the unknown takes values in $\Ecal_\theta$.
Equations \eqref{eq:normal-transport-component} and
\eqref{eq:recursion-normal-compatibility} show that the residual
vector source belongs to $H_\theta$. Its scalar source is
unrestricted. We are therefore left with an inhomogeneous
square system on $\Ecal_\theta$ with invertible derivative
coefficient. Use the uncut frame $\mathcal F$ and the
inhomogeneous part of Lemma~\ref{lem:planar-frames} to solve it
on $D_{R_j}^2\times I_0$. All source terms, including the
particular pair and its derivatives, have transverse support
in $K$; the constructed solution retains that support and
extends smoothly by zero in $t$. This proves the induction.
After finitely many steps, every coefficient is smooth on
$D_{R_N}^2\times\R$, hence on a neighborhood of
$\overline B_{5r}$.

Set
\[
 a^{[N]}=\sum_{j=0}^N\tau^{-j}a_j,
 \qquad b^{[N]}=\sum_{j=0}^N\tau^{-j}b_j.
\]
The first equation of \eqref{eq:physical-recursion} telescopes
to
\begin{equation}
\label{eq:truncated-vector-residual}
 E_\tau^{-1}R(E_\tau a^{[N]},E_\tau b^{[N]})
       =\tau^{-N}R(a_N,b_N).
\end{equation}
The normal constraints telescope separately, giving
\begin{equation}
\label{eq:truncated-divergence}
 E_\tau^{-1}\diverg(E_\tau a^{[N]})
       =b^{[N]}+\tau^{-N}q_N.
\end{equation}
Using the dependence of $R$ on its scalar argument, these
identities give the physical residual explicitly:
\begin{equation}
\label{eq:physical-truncated-residual}
\begin{aligned}
 f_\tau
 &\coloneqq E_\tau^{-1}L(E_\tau a^{[N]})\\
 &=\tau^{-N}\bigl[R(a_N,b_N)+k\nabla q_N
           +(\nabla\lambda)q_N+\tau k\theta q_N\bigr].
\end{aligned}
\end{equation}
All coefficients on the right are fixed smooth functions.
Consequently $\norm{f_\tau}_{L^2(B_{5r})}=O(\tau^{1-N})$.

Corollary~\ref{cor:physical-solvability} gives a correction
$r_\tau\in L^2(B_{5r};\C^3)$ satisfying
\[
 E_\tau^{-1}L(E_\tau r_\tau)=-f_\tau,
 \qquad \norm{r_\tau}_{L^2(B_{5r})}=O(\tau^{1-N}).
\]
It is smooth in $B_{5r}$ by Lemma~\ref{lem:scaled-interior}.
The same Lemma, on $B_{4r}$, gives
\begin{equation}
\label{eq:physical-remainder-two-derivatives}
 \sum_{j=0}^2h^j\norm{\nabla^jr_\tau}_{L^2(B_{4r})}
                      =O(\tau^{1-N}).
\end{equation}
Define its conjugated physical divergence by
\[
 d_\tau=E_\tau^{-1}\diverg(E_\tau r_\tau)
             =\diverg r_\tau+\tau\theta\cdot r_\tau.
\]
It follows directly from
\eqref{eq:physical-remainder-two-derivatives} that
\begin{equation}
\label{eq:physical-divergence-remainder}
 \norm{d_\tau}_{L^2(B_{4r})}
    +h\norm{\nabla d_\tau}_{L^2(B_{4r})}
 \le C\tau\sum_{j=0}^2h^j
                    \norm{\nabla^jr_\tau}_{L^2(B_{4r})}
 =O(\tau^{2-N}).
\end{equation}

Now put $u_\tau=E_\tau(a^{[N]}+r_\tau)$. It is a smooth exact
solution of $Lu_\tau=0$ on $B_{5r}$, and its physical pair
satisfies
\[
 E_\tau^{-1}(u_\tau,\diverg u_\tau)
  =\bigl(a^{[N]}+r_\tau,
         b^{[N]}+\tau^{-N}q_N+d_\tau\bigr).
\]
The smooth finite expansions and the remainder estimates imply
the first assertion of \eqref{eq:physical-leading-limit}.
They also show that, if the last displayed pair is denoted
by $\mathcal A_\tau$, then
\[
 \norm{\mathcal A_\tau}_{L^2(B_{4r})}
       +h\norm{\nabla\mathcal A_\tau}_{L^2(B_{4r})}=O(1).
\]
Since, for $j=1,2,3$,
$E_\tau^{-1}\partial_j(E_\tau\mathcal A_\tau)
=\partial_j\mathcal A_\tau+\tau\theta_j\mathcal A_\tau$,
the second assertion of \eqref{eq:physical-leading-limit}
follows. Repeating the construction for the three initial
columns proves the Proposition.
\end{proof}
